\begin{definition}[Actual closed contractions with inserted bond matrices]%
    \label{def:peps_graph_inserted_bond_state}
    \lean{TNLean.PEPS.graphInsertedBondNetwork}
    \leanok
    \uses{def:peps_graph_bond_contraction}
    Let $\Gamma$ be a finite simple graph with ordered endpoints, let $X$ be
    a finite virtual alphabet, and let $P_v$ be a physical alphabet at each
    vertex. For site coefficients $A_v$ and matrices $K_e$, define
    \begin{align}
        \Psi_K(\sigma)
        =\sum_{\zeta:E(\Gamma)\to X\times X}
        \left(\prod_e(K_e)_{\zeta_{h,e},\zeta_{t,e}}\right)
        \prod_v A_v(\zeta|_v,\sigma_v).
    \end{align}
    Here $\zeta|_v$ reads the head or tail label at the actual incidence.
    Thus each bond matrix has row at the head and column at the tail.
    This is an auxiliary inserted-bond extension of the closed contraction
    calculation in \cite[Section~7]{Schuch2010PEPS}.
\end{definition}

\begin{theorem}[Identity and regular insertions in the actual contraction]%
    \label{thm:peps_graph_inserted_bond_recovery}
    \lean{TNLean.PEPS.graphInsertedBondNetwork_one,
        TNLean.PEPS.graphInsertedBondNetwork_leftRegular_eq_labels,
        TNLean.PEPS.graphInsertedBondNetwork_leftRegular_eq_twisted,
        TNLean.PEPS.graphInsertedBondNetwork_leftRegular_eq_stateCoeff}
    \leanok
    \uses{def:peps_graph_inserted_bond_state}
    When every $K_e=I$, the contraction is the ordinary closed graph
    contraction. For $X=G$ and $K_e=L(u_e)$, where $L$ is the left-regular
    representation, it is exactly the ordinary contraction with the head
    label changed from $\eta_e$ to $u_e\eta_e$ and the tail label retained.
    This equality holds for vertex-dependent physical alphabets. For uniform
    numbered alphabets it is exactly the original twisted tensor-network
    coefficient, with the existing oriented regular insertion convention of
    \cite[Section~6.4]{Schuch2010PEPS}.
\end{theorem}

\begin{proof}\leanok
    The identity matrix constrains the two endpoint labels to agree.
    The regular matrix instead imposes $\zeta_{h,e}=u_e\zeta_{t,e}$.
    Summing the constrained head labels leaves precisely the ordinary
    tail-label sum and the stated site coefficients.
\end{proof}

\begin{theorem}[Weighted averaging contractions with arbitrary bond insertions]%
    \label{thm:peps_graph_inserted_averaging_bond_state}
    \lean{TNLean.PEPS.graphInsertedBondNetwork_dressedAveragingSite,
        TNLean.PEPS.graphInsertedBondNetwork_dressedAveragingSite_group,
        TNLean.PEPS.graphInsertedBondNetwork_averagingSite_group}
    \leanok
    \uses{def:peps_graph_inserted_bond_state,
        thm:peps_monoid_hom_commuting_weight}
    Let $G$ be finite and $U$ a matrix representation on $\mathbb C^X$.
    Dress the averaging site by the same virtual matrix $W$ on each leg.
    Regroup the physical labels into head-tail bond pairs $\beta_e$. Then
    \begin{align}
        \Psi_{U,W,K}(\beta)
        =|G|^{-|V|}\sum_{q:V\to G}
        \prod_e
        [U(q_h)WK_eWU(q_t^{-1})]_{\beta_{h,e},\beta_{t,e}}.
    \end{align}
    This identity holds for arbitrary $W$ and $K_e$.
    If $W$ commutes with every $U(g)$ and $K_e=U(u_e)$, it becomes
    \begin{align}
        \Psi_{U,W,u}(\beta)
        =|G|^{-|V|}\sum_{q:V\to G}
        \prod_e[W^2U(q_hu_eq_t^{-1})]_{\beta_{h,e},\beta_{t,e}}.
    \end{align}
    Taking $W=I$ gives the ordinary averaging-site contraction.
    No connectedness or compatibility among the group insertions is required.
    These finite-simple-graph coefficient identities extend the calculation in
    \cite[Section~7]{Schuch2010PEPS}; the normalized site convention is recorded
    in \cite{gap:rmp_peps_quantum_double_g_isometry}.
    No physical isometry for arbitrary $K_e$ is asserted.
\end{theorem}

\begin{proof}\leanok
    Expand the independent local averages and regroup the incidence products
    by edges. The two independent virtual endpoint sums multiply
    $U(q_h)W$, $K_e$, and $WU(q_t^{-1})$ in that order.
    For representation insertions, commute the weights and combine the
    representation factors to obtain $W^2U(q_hu_eq_t^{-1})$.
\end{proof}
